\section{Justification of MAS}
Although forecasting, imputation, anomaly detection, and classification differ in outputs and objectives, they share a common workflow: (i) \textbf{perceive} salient temporal structures, (ii) \textbf{reason} about global shape constraints and local irregularities, and (iii) \textbf{execute} decisions under validity constraints (e.g., consistency with observations, boundedness, and smoothness). This aligns naturally with an agentic view and motivates modeling time-series analysis as analysis--reasoning--execution pipeline.

A single monolithic model typically entangles these steps into one latent representation with a fixed head, which makes it hard to (a) adapt the execution pathway to heterogeneous series characteristics and task requirements, and (b) incorporate domain tools (e.g., decomposition, robust statistics, constraint projection) in an auditable and controllable manner. A multi-agent design provides a principled way to modularize the workflow into specialized roles while still enabling coordination via a controlled shared-memory interface. Specifically, MAS4TS follows an \textbf{Analyzer--Reasoner--Executor (ARE)} paradigm: the \emph{Analyzer} produces grounded statistical priors and a plot-based interface, the \emph{Reasoner} converts perceived structures into explicit anchors and trajectory constraints, and the \emph{Executor} selects, composes, and verifies task-specific tool chains.

This separation supports adaptive computation (different tool compositions for different regimes), improves interpretability (anchors, constraints, and verifier checks are inspectable), and reduces training burden by focusing learning on coordination-critical components (e.g., routing and memory gating) rather than retraining task-specific backbones. In this sense, the benefit of MAS is not merely parallelization, but the ability to unify diverse time-series tasks via a shared reasoning-and-execution protocol with explicit intermediate artifacts.